\documentclass[11pt]{article}
\usepackage[utf8]{inputenc}
\usepackage[english]{babel}
\usepackage[
backend=biber,
style=alphabetic,
sorting=ynt
]{biblatex}
\addbibresource{references.bib}
\usepackage{amsmath}  % Basic math symbols and environments
\usepackage{amssymb}  % Extended symbol collection
\usepackage{mathtools} % Fixes and additions to amsmath
\usepackage{graphicx}
\usepackage{float}
\usepackage{csvsimple}
\usepackage[colorlinks=true]{hyperref}
\usepackage{indentfirst}
\usepackage{algorithm}
\usepackage{algpseudocode}
\DeclareGraphicsExtensions{.png,.pdf}
\begin{document}
\subsection{Monty Hall favours one door}
	

\begin{align*}
	P(H|D) &= \frac{P(H)P(D|H)}{P(D)} \\
	\\
    H &= C1 = \text{ car behind door \#1} \\
    D &= M{=}2 = \text{ door \#2 chosen} \\
    P(C1|M{=}2) &= \frac{P(C1)P(M{=}2|C1)}{P(D)} \\
    \\
    P(M{=}2|C1) &= 3/4\\
    P(M{=}2|C2) &= 0\\
    P(M{=}2|C3) &= 1\\
    P(M{=}2) &= P(M{=}2|P)*P(C1) + P(M{=}2|C2)*P(C2) + P(M{=}2|C3)*P(C3) \\
    P(M{=}2) &= 3/4*1/3 + 0*1/3 + 1 * 1/3  \\
    	&= 7/12 \\
    \\
    P(C1|M{=}2) &= \frac{1/3 * 3/4}{7/12} \\
    &= \frac{3/12}{7/12} \\
    &= \frac{3}{7}
\end{align*}	

\subsection{Being Dutch Booked}

\begin{align*}
P(prick) &= 0.6 \\
P(!prick) &= 0.6 \\
P(prick or !prick) &= 1\\
\end{align*}

a) This breaks the "sum rule", P(prick) + p(!prick) do not sum to 1. Since there are only two events, and they are disjoint, it also breaks the "addition rule" and the "complement rule".

b) 

\begin{align*}
	\textbf{P(x)} &= 0.6 = \frac{1}{D} = \frac{3}{5} \\
	\textbf{D} &= \frac{5}{3} \\
	\textbf{payout} &= 100 \\
	\textbf{payout} &= \textbf{price} * D \\
	\textbf{price} &= \frac{\textbf{payout}}{D} = \frac{100}{\frac{5}{3}} \\
	\textbf{price} &= 60
\end{align*}

This leads to a standard dutch book. However, it does require bettors on both bets to guarantee a win.

\begin{center}
	\begin{tabular}{|c | c | c | c | c | c |}
		\hline
		\textbf{Bet} & \textbf{P(x)} &  \textbf{Payout} & \textbf{D} &  \textbf{Price} \\
		\hline		
		\textbf{pricked}  & 0.6 & 100 & $\frac{5}{3}$ &  60 \\
		\hline		
		\textbf{!pricked} & 0.6 & 100 & $\frac{5}{3}$ & 60 \\
		\hline
		 & 1.2 & 100 &  & 120 \\
		\hline		 
	\end{tabular}
\end{center}

c) A more subtle
	
\begin{figure}[h]
	\centering
	\includegraphics[width=0.25\textwidth]{assignment1-fig1}
	\caption{undirected graph}
	\label{fig:graph}
\end{figure}

\begin{figure}[h]
	\centering
	\includegraphics[width=0.25\textwidth]{assignment1-digraph}
	\caption{directed graph}
	\label{fig:digraph}
\end{figure}

\end{document}
